% ============================================================
% 个人实验报告（实验三 仿真部分）——桌面物体自动分类整理（mechArm + Gazebo 执行线）
% 编译：xelatex 两遍（中文需 XeLaTeX；依赖 ctex、xeCJK、tikz）
% ============================================================
\documentclass[12pt,a4paper]{ctexart}

\usepackage[margin=2.4cm]{geometry}
\usepackage{amsmath,amssymb}
\usepackage{xcolor}
\usepackage{booktabs}
\usepackage{array}
\usepackage{enumitem}
\usepackage{listings}
\usepackage{hyperref}
\hypersetup{colorlinks=true,linkcolor=blue,urlcolor=blue}
\usepackage{tikz}
\usetikzlibrary{positioning,arrows.meta,shapes.geometric,fit,backgrounds}

% 让 Ⅰ(罗马数字)、≈、≥、≤ 等走 CJK 字体（lmroman 无这些字形）
\xeCJKDeclareCharClass{CJK}{"2160 -> "2160, "2248 -> "2248, "2264 -> "2265}
\emergencystretch=3em
\sloppy

\lstset{
  basicstyle=\ttfamily\small,
  frame=single,
  breaklines=true,
  columns=fullflexible,
  numbers=left,
  numberstyle=\tiny\color{gray},
  aboveskip=4pt,belowskip=2pt,
}

\tikzset{
  proc/.style={draw, rounded corners=2pt, fill=blue!6, align=center, font=\small, inner sep=4pt},
  sys/.style={draw, rounded corners=3pt, fill=green!7, thick, align=center, font=\small, inner sep=5pt},
  seam/.style={draw, rounded corners=2pt, fill=orange!14, thick, dashed, align=center, font=\small, inner sep=4pt},
  hw/.style={draw, rounded corners=3pt, fill=red!7, thick, align=center, font=\small, inner sep=5pt},
  dec/.style={draw, diamond, aspect=2.2, fill=orange!10, align=center, font=\scriptsize, inner sep=2pt},
  arr/.style={-{Stealth[length=2.6mm]}, thick},
  lbl/.style={font=\scriptsize\itshape},
  grp/.style={draw, dashed, rounded corners=4pt, inner sep=6pt},
}

\begin{document}

% ---------------- 封面信息 ----------------
\begin{center}
  {\LARGE\bfseries 个人实验报告（实验三）}\\[6pt]
  {\Large 桌面物体自动分类整理 —— 仿真执行（mechArm + Gazebo，阶段一）}
\end{center}

\begin{center}
\renewcommand{\arraystretch}{1.35}
\begin{tabular}{@{}l p{3.5cm} l p{3.5cm} l p{3.5cm}@{}}
  班级 & \underline{2024届中外2班} & 姓名 & \underline{马晨曦} & 学号 & \underline{24020036047} \\
  小组角色 & \multicolumn{3}{l}{任务控制与系统集成 · \textbf{仿真执行主责}} & 日期 & \underline{2026-09-24} \\
\end{tabular}
\end{center}

\section{实验的基本要求}

本实验为机器人集成小组项目Ⅰ《桌面物体自动分类整理》，要求在实验一（视觉检测模型）与实验二
（机械臂定点抓取）的基础上，做出一个\textbf{“看得见—判得准—抓得起—放得对”的全自动整理系统}：
相机看桌面，识别出网格里摆放的多类物体，机器人依次把每个物体按类别放进对应料盒，直到桌面清空。

实验要求分为\textbf{阶段一 仿真开发}与\textbf{阶段二 真机验证}两段，且
\textbf{阶段二的准入条件 = 阶段一仿真验收通过}（顺序不可颠倒）。
\textbf{本人承担的是阶段一仿真线}——mechArm 270 机械臂 + Gazebo + ROS\,2 Humble，
负责场景搭建与联调、检测链路接入、取放执行链、仿真演示视频，以及仿真与真机的尺寸对齐。
\textbf{本报告围绕这条线展开。}

\textbf{本版范围声明：}本报告\textbf{只写本人承担的仿真线}。阶段二真机（RoboMaster EP）
由组内其他同学承担，本版不写真机的标定数据与运行结论，也不代填任何估计值。
两条线唯一共享的是\textbf{任务控制核心 \texttt{sort\_core}、接口契约 \texttt{PickPlace}
与原因枚举 / 日志 schema}——这部分由本人建立、两线共用（见 §3.1），
因此本报告在讲它们时说的是\textbf{这套核心本身}，与跑在哪条线上无关。

\subsection*{阶段一仿真的验收结果（本人的主线成果）}

仿真线（mechArm + Gazebo，ROS\,2）的\textbf{记分轮}已于 2026-09-24 跑通，
\begin{center}
\begin{tabular}{@{}ccccc@{}}
\toprule
格号 & 类别 & 目标料盒 & 落料槽位 & 结果 \\
\midrule
c1 & cup   & bin\_cup   & 0 & 成功 \\
c2 & mouse & bin\_mouse & 0 & 成功 \\
c3 & cup   & bin\_cup   & 1 & 成功 \\
c4 & cup   & bin\_cup   & 2（第 2 层） & 成功 \\
c5 & mouse & bin\_mouse & 1（叠层） & 成功 \\
c6 & cup   & bin\_cup   & 3 & 失败（放置偏差 52/87/17\,mm） \\
\bottomrule
\end{tabular}
\end{center}
结果文件 \texttt{exp3\_logs/run\_20260924\_112831/result.json} 给出
\texttt{verdict=PASS}：场上物体 6 个，正确放入对应料盒 5 个（及格线 5），
\textbf{规划失败 0 次}。

\textbf{异常轮}（\texttt{run\_20260924\_114058}）单独跑：c3 一格故意空置，
\texttt{verdict=TRIAL}。空网格被跳过的证据不是某一行写了什么，而是
\texttt{records.jsonl} 里\textbf{没有 c3 这一行}、\texttt{result.json} 的
\texttt{reasons} 里也没有 c3 的键；其余 5 个全部放置成功。这一轮必然是 TRIAL 而非
PASS —— 判定规则只对\textbf{完整}的验收轮下结论，所以异常轮必须与记分轮分开跑。

\textbf{口径说明（必须分开陈述）。}本次仿真跑在没有 GPU 的云机容器里，Gazebo 走软件渲染，
跑到第 3$\sim$4 轮会出现\textbf{画面停滞}：帧戳照常推进而画面内容一字不变，依赖图像的
HSV 检测因此失效。为把故障隔离在渲染这一段、继续验证下游链路，记分轮改用\textbf{真值投影}
生成检测框：物块世界坐标经同一份已标定的相机参数投影为像素，再走完全相同的
“像素 $\to$ 桌面 $\to$ 格子”映射，相机标定与映射链路未作任何改动。因此这一轮
\textbf{能证明}的是相机标定、投影与网格映射正确，以及规划—取放—落料—回航这条执行链可行；
\textbf{不能}声称 HSV 检测算法在\textbf{整轮}任务中稳定有效——渲染在第 3$\sim$4 轮停滞之后，
就再也拿不到全程证据了。\textbf{可以单独保留的是：}HSV 在渲染健康的\textbf{首轮}把 6 个物体
全部正确映射（逐 blob 的“类别@像素$\to$格子”记录），以及它实际处理到的格子类别与摆放表
一致。\textit{（该 HSV 轮的原始日志未随证据目录回收，此条以《实验3\_仿真验收通过\_20260924.md》
第二节的逐轮表为据；真值检测器那一轮的 6/6 另有一份可复核的日志。）}

\textbf{遗留：}叠第 2 层时物块之间发生互穿，LCP 求解器迭代数冲到 122847 并把已放好的物块弹飞；
第 2 层落差 \texttt{dz}=0.036\,m 仅比杯高 0.035\,m 多 1\,mm，加大落差或让层间 xy 错开是下一步
要改的地方。c6 是记分轮唯一失败项。

\subsection*{仿真任务（docx §五）}

\begin{itemize}[leftmargin=1.6em]
  \item 在 Gazebo 场景中搭建\textbf{取物网格（2$\times$3 共 6 格）与 2 个分类料盒}，
        摆放 \textbf{2 类共 6 个物体}（4 个 cup + 2 个 mouse）；
  \item 通过\textbf{检测链路}识别物体类别与所在网格——HSV \texttt{color\_detector}
        或真值投影 \texttt{truth\_detector}（§3.4）；
  \item 用\textbf{单条 Launch} 启动相机、识别、机械臂与任务控制，把不同类别放入对应料盒。
\end{itemize}

\subsection*{验收标准（docx §六）}

\begin{itemize}[leftmargin=1.6em]
  \item 6 个物体中\textbf{至少 5 个被正确分类整理}；
  \item \textbf{正常任务过程中不得进行人工类别判断和动作触发}；
  \item 正确处理\textbf{至少 2 类异常}（空网格、未识别物体应被跳过；抓取失败或目标不可达时
        应返回安全位置或停止）；
  \item 提供一个入口启动图像、识别、机械臂和任务控制；
  \item 全程保存\textbf{识别、抓取和任务状态日志}。
\end{itemize}

\subsection*{仿真平台的确定}

docx §实施原则写的是“真实相机、Jetson 和 mechArm”。阶段一\textbf{正是这条款的原生形态}：

\begin{itemize}[leftmargin=1.6em]
  \item \textbf{机械臂}：mechArm 270（6 自由度），由 \texttt{mecharm\_grasp} 接入 Gazebo，
        走 \texttt{joint\_trajectory\_controller}；逆解用 \texttt{c4\_2/core/planner.py}
        的 Newton + 数值 Jacobian——关节限位里 \texttt{q5} 上限 2.0071\,rad
        是搬运段最常撞到的一条（§4 问题 2）；
  \item \textbf{相机}：Gazebo 的 \texttt{overhead\_camera}（俯瞰，640$\times$480，10\,Hz），
        经\textbf{相机标定}把像素映射到桌面坐标系（\texttt{task.json} $\to$
        \texttt{camera.calibration}）；
  \item \textbf{运行环境}：AutoDL 云机容器，\textbf{无 GPU}，Mesa 走 llvmpipe 软件渲染。
        这一点后面直接决定了检测链路的选型（§3.4）。
\end{itemize}

\textbf{注：}本人在这条线上承担\textbf{任务控制（状态机）/ 接口契约 / 日志与验收口径}
这条主线，并负责\textbf{仿真场景搭建与联调、取放执行链、仿真演示视频}，
以及\textbf{仿真与真机的尺寸对齐}。检测模型来自实验一的公开仓库；
阶段二真机（RoboMaster EP）的执行侧由组内其他同学承担。

\section{个人分工}

本人在阶段一仿真线负责以下五件事：

\begin{enumerate}[leftmargin=1.8em]
  \item \textbf{仿真场景与单 Launch}：Gazebo world 里搭出 2$\times$3 取物网格、2 个分类料盒
        与俯瞰相机，并把相机、识别、机械臂、任务控制\textbf{收进一条
        \texttt{ros2 launch exp3\_sim exp3\_sim.launch.py}}（docx §六 最后一条要求）；
  \item \textbf{取放执行链 \texttt{GazeboFollowExecutor}}：实现执行契约的\textbf{六段动作}
        （下探/夹取/抬离/搬运/松放/退回），并按段号映射为动作级原因；其中
        \textbf{同料盒多物错位落料}与\textbf{搬运高度阶梯}是让 6 个物体能全部放下的关键（§4）；
  \item \textbf{检测链路接入与真值投影检测器}：在容器无 GPU、Gazebo 软渲染会停帧的条件下，
        把检测框的生成换一种来源而\textbf{不动相机标定与映射链路}（§3.4）；
  \item \textbf{任务控制核心 \texttt{sort\_core} / 接口契约 \texttt{PickPlace} /
        原因枚举与日志 schema}：这套核心是本人建立的，\textbf{两条线共用同一份}，
        真机线只是把执行侧换成 EP，核心一行未改（见 §3.1）；
  \item \textbf{仿真运行与取证}：三轮实跑（记分轮 / 异常轮 / 录屏轮）、演示视频录制，
        以及证据目录与验收日志的整理。
\end{enumerate}

组内分工概览：\textbf{本人} = 仿真场景搭建与联调 + 取放执行链 + 仿真演示视频
+ 仿真与真机的尺寸对齐 + 任务控制与契约；\textbf{组内其他同学} = 阶段二真机（RoboMaster EP）
执行侧、识别链路（真模型）与真机相机标定。

\section{项目实现的原理}

\subsection{双线架构：一套任务核心，两个注入缝隙}

任务控制核心 \texttt{TaskController} 只依赖两个函数缝隙 \texttt{scan()}（报一帧桌上检测）
与 \texttt{pick\_place(cell\_id, cls)}（把指定格子的物体取走放好），
\textbf{自己不 import 任何 ROS 类型}。仿真线要做的只是\textbf{填好这两个缝隙的实现}
（图~\ref{fig:arch}）：

\begin{itemize}[leftmargin=1.6em]
  \item 缝隙 1：订阅 \texttt{/detections\_d2a}（\texttt{Detection2DArray}），
        把检测框中心经相机标定映射成桌面坐标与格号；
  \item 缝隙 2：发 \texttt{PickPlace} Action，由 \texttt{GazeboFollowExecutor}
        完成六段动作并回传动作级原因。
\end{itemize}

这条“决策与通信解耦”的设计是本人定的，\textbf{阶段一仿真先把它跑通}；正因为缝隙两侧口径固定，
后来换到真机时\textbf{只换了这两个实现，契约、原因枚举、日志 schema 一行未改}。

\begin{figure}[htbp]
  \centering
  \begin{tikzpicture}[node distance=3.4mm]
    \node[sys, text width=11.4cm] (cam) {\textbf{Gazebo 俯瞰相机 \texttt{overhead\_camera}}
      （640$\times$480，10\,Hz）\\ $\rightarrow$ \textbf{检测链路}（HSV 或真值投影）
      $\rightarrow$ \texttt{/detections\_d2a}};
    \node[seam, text width=11.4cm, below=of cam] (gap1) {\textbf{缝隙 1：\texttt{scan()}}
      —— 检测框中心 $\to$ 像素 $\to$ 桌面坐标 $\to$ 格号};
    \node[sys, text width=11.4cm, below=of gap1] (ctl) {\textbf{任务控制核心 \texttt{TaskController}}
      （纯 Python，无 ROS 依赖）\\ 扫描 $\rightarrow$ 逐条判定 $\rightarrow$ 选目标 $\rightarrow$
      执行 $\rightarrow$ 循环，直到清空 / 异常退出};
    \node[seam, text width=11.4cm, below=of ctl] (gap2) {\textbf{缝隙 2：\texttt{pick\_place(cell\_id, cls)}}
      —— \texttt{GazeboFollowExecutor}（goal 只给格号与类别）};
    \node[sys, fill=blue!8, text width=11.4cm, below=of gap2] (exe) {\textbf{\texttt{GazeboFollowExecutor}}
      ：把格号/类别解成\textbf{桌面坐标 + 料盒槽位}，执行六段动作 $\rightarrow$
      \textbf{关节轨迹 + 夹爪控制器}};
    \node[hw, text width=11.4cm, below=of exe] (rm) {\textbf{mechArm 270 + Gazebo}：6 自由度臂
      + 夹爪\quad$\vert$\quad\textbf{ROS\,2 Humble}（\texttt{joint\_trajectory\_controller}）};
    \node[sys, fill=green!10, text width=11.4cm, below=of rm] (log) {\textbf{日志/验收}：
      \texttt{exp3\_logs/run\_*/}\ \texttt{records.jsonl} $+$ \texttt{result.json}
      $+$ \texttt{run.log}};

    \draw[arr] (cam)--(gap1); \draw[arr] (gap1)--(ctl); \draw[arr] (ctl)--(gap2);
    \draw[arr] (gap2)--(exe); \draw[arr] (exe)--(rm); \draw[arr] (rm)--(log);
    \node[lbl, right=1mm of gap1, text=orange!60!black, align=left] {仿真：订阅话题};
    \node[lbl, right=1mm of gap2, text=orange!60!black, align=left] {仿真：发 Action};
  \end{tikzpicture}
  \caption{仿真线主链与两个注入缝隙（橙色虚线）。缝两侧口径固定，故任务核心与执行侧可分别开发、
  失败可分别定位；这两处缝隙正是后来换到真机时唯一被替换的地方。}
  \label{fig:arch}
\end{figure}

\subsection{一次取放拆成六段：把动作序列固定，把“在哪儿”参数化}

“取一个物体并放好”是一个持续数秒的长任务，服务端把它切成固定的\textbf{六个阶段}，
每进入一段先发进度，再交给执行者 \texttt{PickExecutor}。
这个抽象是相对实验二最重要的一步泛化：实验二的直驱程序是“一个固定取物点 $\to$ 固定放置点”
的循环，而本实验有 6 格 $\times$ 2 料盒；若直接套用就会退化成“每格写死一个周期”。
真正会变的只是“\textbf{目标和料盒在哪儿}”，动作序列本身完全相同 ——
于是把序列固定成六个方法、把位置压进参数（表~\ref{tab:stage}）。

\begin{table}[htbp]
  \centering
  \small
  \begin{tabular}{@{}c l p{5.4cm} p{4.4cm}@{}}
    \toprule
    段 & 语义 & 仿真侧实现 & 该段失败时的原因 \\
    \midrule
    1 & 下探 & 规划到格子上方 $\to$ 下降到抓取高度 & \texttt{unreachable\_target} \\
    2 & 夹取 & 夹爪控制器闭合 & \texttt{exec\_error} \\
    3 & 抬离 & 抬到搬运高度 \texttt{carry\_z} & \texttt{grasp\_failed} \\
    4 & 搬运 & 横移到目标料盒\textbf{槽位}上方（高度按阶梯由高到低退让） & \texttt{unreachable\_target} \\
    5 & 松放 & 下降到放置高度 $\to$ 开爪 & \texttt{dropped\_or\_}\allowbreak\texttt{missed\_place} \\
    6 & 退回 & 抬离 $\to$ 回停靠位（\texttt{DOCK\_J1}） & \texttt{unreachable\_target} \\
    \bottomrule
  \end{tabular}
  \caption{六段动作在仿真侧的实现，以及每段失败时映射出的动作级原因。}
  \label{tab:stage}
\end{table}

这张表背后有一条我给自己定的规则，值得单独说：\textbf{只报这一段的代码真能做出来的判断，
不替谁宣布一个没人做过的判定。} 两个具体后果：

\begin{itemize}[leftmargin=1.6em]
  \item \textbf{夹取段报 \texttt{exec\_error} 而不是 \texttt{grasp\_failed}}：这一段的代码只是
        合一下夹爪，\textbf{从没判过“夹住没有”}（判定在下一段抬离）；而
        \texttt{grasp\_failed} 在原因枚举里的定义恰恰是“抬离后判定没夹起”。填它等于
        \textbf{凭空宣布一个没做过的判定}——查日志的人会以为系统真的判过。
  \item \textbf{搬运段与退回段报 \texttt{unreachable\_target}}：仿真侧有\textbf{规划器}，
        “够不着”是它给出的真实结论（逆解不收敛或撞上关节限位，如 \texttt{q5} 上限
        2.0071\,rad）。真机线没有规划层、报的是驱动硬错 \texttt{exec\_error}——
        \textbf{这不是没对齐，而是两个平台上“发生了什么”本来就不是同一件事。}
\end{itemize}

\begin{figure}[htbp]
  \centering
  \begin{tikzpicture}[node distance=2mm]
    \node[proc, fill=blue!10] (s1) {\scriptsize 1\\下探\\\texttt{goto+moveto}};
    \node[proc, fill=blue!10, right=of s1] (s2) {\scriptsize 2\\夹取\\\texttt{close}};
    \node[proc, fill=blue!10, right=of s2] (s3) {\scriptsize 3\\抬离\\\texttt{moveto}$+$判};
    \node[proc, fill=blue!10, right=of s3] (s4) {\scriptsize 4\\搬运\\\texttt{goto}};
    \node[proc, fill=blue!10, right=of s4] (s5) {\scriptsize 5\\松放\\\texttt{open}$+$判};
    \node[proc, fill=blue!10, right=of s5] (s6) {\scriptsize 6\\退回\\\texttt{moveto+goto}};
    \node[hw, fill=red!10, below=7mm of s3] (f) {\scriptsize 任一段失败 $\Rightarrow$
      \textbf{立刻收尾}：开爪 $\to$ 抬到高 $\to$ 回停靠位};
    \draw[arr] (s1)--(s2); \draw[arr] (s2)--(s3); \draw[arr] (s3)--(s4);
    \draw[arr] (s4)--(s5); \draw[arr] (s5)--(s6);
    \draw[arr, dashed, red!70] (f.north) -- (s3.south);
  \end{tikzpicture}
  \caption{六段动作链（仿真侧）。同一套段号也被真机线沿用，差别只在各段的原语是什么。
  失败段号 $\to$ 原因 $\to$ 收尾这条路径见 §3.5。}
  \label{fig:stages}
\end{figure}

\subsection{格号怎么变成机械臂能用的坐标：相机标定与桌面坐标系}

“格号 $\to$ 坐标”这一步在仿真侧用的是\textbf{桌面坐标系}，它由\textbf{相机标定}定义：

\begin{itemize}[leftmargin=1.6em]
  \item \textbf{相机标定}（\texttt{task.json} $\to$ \texttt{camera.calibration}）给出
        \texttt{offset\_px} 与 \texttt{scale\_px\_per\_m}，把\textbf{像素}线性映射到
        \textbf{桌面坐标}（米）。检测框中心先$\to$像素$\to$桌面坐标$\to$格号，链路全程可查。
  \item \textbf{格号与料盒}各自在 \texttt{grid\_cells.json} / \texttt{bins.json} 里给出桌面坐标，
        执行侧\textbf{只认坐标不认格号}——换格子、换料盒、挪相机都只改配置，
        \textbf{不动代码、不动接口}。
  \item \textbf{代价（如实记录）：}这套映射是\textbf{线性近似}（缩放 + 平移）。
        它在俯瞰取景下够用；若相机明显倾斜、透视不可忽略，就需要换成单应变换——
        本线沿用线性近似的理由是\textbf{当前取景下的误差可接受}，这一点有
        “6 格全部正确映射”的实跑结果作证（§1）。
\end{itemize}

料盒 \texttt{bins.json} 除了中心坐标，还带\textbf{落料槽位 \texttt{slots}}（每个槽位一组
\texttt{dx/dy/dz} 偏移）与\textbf{按料盒设定的末端偏航 \texttt{yaw\_deg}}——
这两项是后面解决“同一料盒放多个物体”的关键（§4 问题 3）。

\subsection{检测框从哪来：HSV 与真值投影共用同一份标定（本线最关键的一处取舍）}

\subsubsection*{问题的由来}

仿真跑在\textbf{没有 GPU}的云机容器里，Gazebo 走 Mesa 软件渲染（llvmpipe）。
用 HSV \texttt{color\_detector} 跑整轮任务时，实测它吃 \textbf{479\% CPU}（约 5 核），
\texttt{gzserver} 吃到 1013\%；跑到第 3$\sim$4 轮就出现\textbf{画面停滞}——
帧戳照常推进、画面内容却一字不变，依赖图像的检测随之失效。

\subsubsection*{做法：只换“框从哪来”，不动标定与映射}

\begin{enumerate}[leftmargin=1.8em]
  \item 新增 \texttt{truth\_detector}：从 \texttt{/gazebo/model\_states} 取物块世界坐标，
        经\textbf{同一份已标定的相机参数}（§3.3）投影为像素，再按\textbf{归一化}约定
        发出 \texttt{Detection2DArray}；
  \item 下游“像素 $\to$ 桌面 $\to$ 格号”的映射\textbf{一行未改}，
        launch 用 \texttt{detector:=hsv|truth} 切换；
  \item CPU 从 479\% 降到 \textbf{4.5\%}，渲染随之恢复健康——录屏轮抽 7 个时间点的帧校验，
        md5 两两全不同，\textbf{画面全程在变}。
\end{enumerate}

\subsubsection*{为什么必须把话说清楚}

真值投影是\textbf{理想感知}：无漏检、无噪声、无遮挡，比任何真实相机都强。因此：

\begin{itemize}[leftmargin=1.6em]
  \item \textbf{能证明}：相机标定正确、投影与网格映射正确、规划—取放—落料—回航这条执行链可行；
  \item \textbf{不能声称}：HSV 检测算法在\textbf{整轮}任务中稳定有效——渲染停帧之后拿不到
        全程证据。能作为“检测算法有效”证据的只有 HSV 在渲染健康\textbf{首轮}的 6/6 正确映射，
        这一份与真值方案的 6/6 \textbf{是两回事}，不能互相顶替：前者证明的是 HSV 本身，
        后者证明的是标定与执行链。
\end{itemize}

这两件事在报告里必须分开陈述。混为一谈，等于把“执行链跑通”说成了“检测算法全程有效”，
或者反过来把理想感知的成绩算到了 HSV 头上。

\subsection{失败即收尾：段号 \texorpdfstring{$\to$}{->} 原因 \texorpdfstring{$\to$}{->} 立刻回到已知状态}

一次取放是\textbf{持续数秒的长动作}。中途失败时机械臂停在一个\textbf{未定义的位置}：
若只记一笔就继续下一个物体，下一次规划会\textbf{从一个非法起点出发}——
轻则逆解不收敛，重则拖着臂横扫桌面。本线的处理是三步（图~\ref{fig:fail}）：

\begin{enumerate}[leftmargin=1.8em]
  \item \textbf{段号定因：}失败发生在第几段，就决定了动作级原因是哪个（表~\ref{tab:stage}），
        日志里能直接看出“是没夹起”还是“够不着”；
  \item \textbf{立刻收尾：}开爪 $\to$ 抬到高点 $\to$ 回停靠位（\texttt{DOCK\_J1}），
        \textbf{尽力而为、不再抛错}（反正已经在失败路径上）；
  \item \textbf{每次扫描前强制回停靠位：}臂若停在桌面中间会\textbf{挡住俯瞰相机}，
        实测直接造成漏检。所以收尾之后强制回停靠位，并在两次扫描之间设
        \textbf{最小间隔}（8\,s），等画面稳定下来再取帧。
\end{enumerate}

\begin{figure}[htbp]
  \centering
  \begin{tikzpicture}[node distance=4mm and 7mm]
    \node[proc] (base) {六段执行中};
    \node[dec, right=of base] (q) {某段失败？};
    \node[proc, right=of q, fill=orange!12] (map) {按\textbf{段号}定因\\（动作级原因）};
    \node[hw, right=of map, fill=red!8] (rst) {\textbf{立刻收尾}\\开爪 $\to$ 抬到高\\$\to$ 回停靠位};
    \node[proc, below=6mm of base, fill=green!10] (nx) {下一轮扫描：\\先回停靠位，再等画面稳定};
    \draw[arr] (base)--(q);
    \draw[arr] (q)--node[above,lbl]{是}(map);
    \draw[arr] (map)--(rst);
    \draw[arr] (rst.south) |- (nx.east);
    \draw[arr] (q.south) |- (nx.east);
  \end{tikzpicture}
  \caption{失败路径：段号定因 $\to$ 立刻收尾 $\to$ 下一次动作前再归位。
  三级递进都是为了同一件事——\textbf{不让机械臂带着未知位姿继续工作}。}
  \label{fig:fail}
\end{figure}

\subsection{日志与验收口径：让“这一轮算不算”由文件说了算}

仿真的日志落在 \texttt{exp3\_logs/run\_<时间戳>/}：\texttt{records.jsonl} 是
\textbf{逐物体}记录（格号、类别、结果、原因），\texttt{result.json} 是\textbf{整轮}结论，
\texttt{run.log} 是过程流水。前两份的 schema \textbf{与真机线完全相同}，
两线结果可以直接比对。

\texttt{result.json} 里定胜负的是这几个字段：

\begin{itemize}[leftmargin=1.6em]
  \item \texttt{verdict}：\texttt{PASS} / \texttt{FAIL} / \texttt{TRIAL}——
        \textbf{只有完整的验收轮才可能判 PASS/FAIL}；场上物体不足时判 \texttt{TRIAL}；
  \item \texttt{placed\_ok} / \texttt{objects\_seen}：“放对了几个”与“场上共几个”；
  \item \texttt{reasons}：按原因码计数（\texttt{out\_of\_grid}、
        \texttt{dropped\_or\_missed\_place}、\texttt{no\_exec} 等）。
\end{itemize}

这里有一条我特意定死的口径：\textbf{空网格的判据是“没有记录”，而不是“记了一条跳过”。}
异常轮把 c3 故意空置之后，\texttt{records.jsonl} 里\textbf{根本没有 c3 这一行}，
\texttt{reasons} 里也没有 c3 的键。这样“跳过了”这件事\textbf{无法靠事后补一行来伪造}——
证据是那份文件的\textbf{缺席}，不是某一行的措辞。

仿真侧还有一点比真机干净：\textbf{整轮没有任何人工参与}。判类别由检测链路给、
选目标与触发动作由状态机给，日志里不存在需要人回答的环节，
因此也不存在“这份结果够不够得上验收凭据”的争议。

\section{项目中遇到的问题和解决方法}

\textbf{问题 1：搬运高度一抬高，规划就失败——而它和搬运距离无关。}

\textit{定位：}搬着物体横移时要抬到一个“既够高不碰桌、又够低够得着”的高度。
一开始按料盒高度的比例取 \texttt{carry\_z} = 0.125 / 0.112\,m，结果 cup 那 4 格
\textbf{全部不可达}。直觉上“离得越远就抬得越高”，但逆解撞的是\textbf{关节限位}
（\texttt{q5} 上限 2.0071\,rad），与目标料盒在 $y$ 方向的\textbf{位置}有关，
与搬运距离无关——这个反直觉的结论只在离线扫参时才看得出来。

\textit{解决：}把可达性判断\textbf{从云机搬到本地}。\texttt{c4\_2/core/planner.py} 的
\texttt{plan\_segment} 只是 Newton + 数值 Jacobian，\textbf{不碰 ROS、不碰 Gazebo}，
所以能离线扫。扫出的上界是：目标在 $-y$ 侧（cup 料盒）时 \texttt{carry\_z} 上限
0.096\,m，$+y$ 侧（mouse 料盒）上限 0.112\,m。于是把阶梯改成\textbf{由高到低退让}
（0.112 $\to$ 0.096 $\to$ 0.088 $\to$ 0.073），并且\textbf{离边界留约 8\,mm 余量}：
瞄准边界时跟踪过冲 1\,mm，就会让下一步规划从一个非法起点出发——
那样成功是偶然、失败是必然。

\textit{个人思考：}“能算出来”和“能稳定跑出来”是两件事。逆解在边界上数学可解，
但物理侧有跟踪误差。由此定了一条规矩：\textbf{凡能被离线验证的判据，就抽到不依赖
运行环境的模块里}（后来摆放表、异常轮空格、高度阶梯都抽到了 \texttt{scene.py}，
因而进得了离线单测）；需要连 Gazebo 才能验证的东西，一轮要几分钟，不值得用来扫参数。

\medskip
\textbf{问题 2：一个料盒要放 3$\sim$4 个物体，第 2 个下降时直接穿进第 1 个里。}

\textit{定位：}6 个物体只分 2 类，一个盒子要放 3$\sim$4 个。都瞄准料盒中心时，
第 2 个物体在\textbf{下降途中}就穿进第 1 个的体内，ODE 的 LCP 求解器爆解，
把已经放好的那个\textbf{炸飞}——实测 c1 被抛到 $(-75.2,\,47.7)$\,m，整轮就此报废。

\textit{解决：}给每个料盒加\textbf{落料槽位 \texttt{slots}}（每个槽位一组
\texttt{dx/dy/dz} 偏移，第二层再叠一个 \texttt{dz}），并按料盒注入\textbf{末端偏航
\texttt{yaw\_deg}}（bin\_cup 取 $-45^\circ$）——料盒 $+x$ 半边在 $\theta=0$ 时不可达，
转过偏航后两个槽位全通。改完后这一项的 LCP 迭代数从 3263 掉到 \textbf{0}。

\textit{个人思考：}“放进去”在仿真里不是几何问题而是\textbf{接触动力学}问题：
物体之间不互穿不是渲染效果，是求解器能收敛的前提。而求解器一旦爆解，
症状（物体乱飞、坐标离谱）与原因（落料位置重合）隔了十万八千里。
\textbf{遗留：}第二层落差 \texttt{dz}=0.036\,m 只比杯高 0.035\,m 多 1\,mm，
加上放置误差仍会互穿（记分轮 LCP 冲到 122847），加大落差或让层间 xy 错开是下一步。

\medskip
\textbf{问题 3：Gazebo 画面停住了，但帧戳还在往前走。}

\textit{定位：}用 HSV 检测器跑整轮任务时观察到\textbf{画面停滞}——帧戳照常推进，
画面内容却一字不变；而同一时刻 \texttt{trace\_blocks.log} 里物块的世界坐标仍在持续更新，
臂的位姿也是实时的。\textbf{物理侧完全正常}，所以死的不是算法而是\textbf{渲染管线}。
根因指向容器无 GPU、Mesa 走 llvmpipe 软件渲染时的 CPU 争抢：
\texttt{color\_detector} 实测吃 \textbf{479\%} CPU（约 5 核）、\texttt{gzserver}
\textbf{1013\%}，渲染抢不到资源。

\textit{（如实标注证据等级：这一轮是 HSV 的停帧轮，其\textbf{原始日志未随证据目录回收}，
上面这些数字以《实验3\_仿真验收通过\_20260924.md》第二节的逐轮表为据；
可独立于该文档复核的是换检测器之后的抽帧校验——见下。）}

\textit{解决与口径：}见 §3.4——改用真值投影给检测框（CPU 降到 4.5\%），
并如实区分“证明了什么、不能证明什么”。

\textit{个人思考：}\textbf{只比对检测数量永远看不出来。}真正一眼定案的是给扫描日志加上
每个 blob 的“类别@像素$\to$格子”之后：同一轮里 c1 已经被取走，却仍在原像素位置被检出
⇒ 画面是\textbf{滞后的旧内容}。诊断日志要打\textbf{可比对的具体值}，不要只打计数。

\medskip
\textbf{问题 4：扫描时机械臂就横在相机视野里，白丢一半物体。}

\textit{定位：}俯瞰相机装在正上方，机械臂的停靠位若还在桌面中间，扫描时它就把整桌挡掉一块。
表现是“明明摆着 6 个物体，只检出 4 个”，而检测算法本身没有问题——
于是很容易往“阈值没调好”的方向查，越查越远。

\textit{解决：}两件事必须配套做：① 把\textbf{停靠位移到画面边缘}
（\texttt{DOCK\_J1} = $-2.6$），并把 \texttt{motion['home']} 指向停靠位；
② \textbf{收尾之后强制回停靠位}，且两次扫描之间设\textbf{最小间隔}（8\,s），
等画面稳定下来再取帧。少做其中一件，就会分别退化成
“遮挡漏检 / 每轮都遮挡 / 取到上一帧的旧内容”。

\textit{个人思考：}感知与执行共用同一个物理空间时，“机器人自己挡住自己的眼睛”
是最容易被忽略的一类耦合。它不报错，只让成功率悄悄下降——
\textbf{没有“把臂挪开”这一步，前面所有调参都是在噪声里调。}

\medskip
\textbf{问题 5：代码明明改对了，收到的却还是旧格式——因为上一轮的进程还活着。}

\textit{定位：}真值检测器与 HSV 检测器发的是\textbf{同一个话题} \texttt{/detections\_d2a}。
切换检测器时若只按新代码重启，上一轮那个\textbf{旧代码的进程}可能仍在发帧，两路数据混在一起。
现象极具误导性：归一化约定已经改对了，日志里收到的却还是像素值 273.0。

\textit{解决：}把\textbf{每一个会往同一话题发帧的进程名}都写进 \texttt{reset\_scene.sh}
的清理名单，并且\textbf{连杀多轮}（launch 会拉一串子进程，父进程被杀后子进程仍会短暂存活，
它们的定时器还会去连新 gzserver 并提前激活控制器）。同类地，用 \texttt{setsid} 起的
记录脚本（轨迹 / 相机 / 时钟 / 录屏）\textbf{不在 launch 的进程树里}，
也必须按名字单独清——漏掉之后它们会一轮轮堆积空转，实测三轮叠加就把 load 顶了上去。

\textit{个人思考：}“我改的代码没生效”，有一半不是代码问题而是\textbf{有两个进程在跑}。
凡存在“两个进程可能发同一路数据”的场景，清理名单就该与代码一起维护。

\subsection*{附：本人这条线的复现方式}

\begin{lstlisting}[numbers=none]
# 1) 离线全套单测（无需 ROS、无需 Gazebo、无需相机）
cd exp3 && <项目 venv>/bin/python -m unittest discover -s tests -t .
#   260 项全绿。注意必须用项目 venv：换解释器会因缺 numpy 静默跳过 25 项

# 2) 同步权威源码 + 重建（图元化 URDF + colcon build）
cd 实验3_仿真运行证据_20260923 && bash resume.sh

# 3) 记分轮（真值检测器）
ssh autodl-exp3 'cd /root/autodl-tmp/ws_mecharm && DET=truth bash runf2.sh'

# 4) 异常轮（c3 故意空置；verdict 必为 TRIAL，与记分轮分开跑）
ssh autodl-exp3 'cd /root/autodl-tmp/ws_mecharm && \
  DET=truth LAUNCH_EXTRA="empty_cell:=c3" bash runf2.sh'

# 5) 录演示视频（订阅 /overhead/image_raw 写 mp4，不依赖 GL）
python3 cam_record.py exp3_sim_raw.mp4 --seconds 400

# 6) 清理（务必在下一轮之前跑；漏了会有旧进程继续发帧）
ssh autodl-exp3 'cd /root/autodl-tmp/ws_mecharm && bash reset_scene.sh'

# 日志：exp3_logs/run_<时间戳>/{run.log, records.jsonl, result.json}
\end{lstlisting}

\vspace{2mm}
{\small 说明：上列第 1 项（260 个离线用例，含摆放表、落料槽位、异常轮空格、
原因映射与跳过语义）在当前仓库中可直接执行并已全绿；\textbf{第 2--6 项需要有
Gazebo 与 ROS\,2 的运行环境}，本报告引用的三轮结果均在其中跑出，证据落在
\texttt{实验3\_仿真运行证据\_20260923/}。相关代码在
\texttt{exp3/ros2/exp3\_sim/}（场景、取放服务端、真值检测器、launch）与
\texttt{exp3/config/}（网格、料盒与槽位、相机标定）中。}
\end{document}
